\documentclass[10pt]{article}

\usepackage[margin=0.58in]{geometry}
\usepackage[default]{sourcesanspro}
\usepackage[T1]{fontenc}
\usepackage[hidelinks]{hyperref}
\usepackage{enumitem}
\usepackage{titlesec}
\usepackage{microtype}

\pagestyle{empty}
\setlength{\parindent}{0pt}
\setlength{\parskip}{0pt}

\setlist[itemize]{
    leftmargin=1.15em,
    topsep=2pt,
    itemsep=2.5pt,
    parsep=0pt,
    partopsep=0pt
}

\titleformat{\section}
  {\large\bfseries}
  {}
  {0pt}
  {}
  [\vspace{-5pt}\rule{\textwidth}{0.4pt}]

\titlespacing*{\section}{0pt}{7pt}{4pt}

\begin{document}

% ============================================================
% HEADER
% ============================================================

\begin{center}
    {\LARGE\textbf{Oth\'on Gonz\'alez}}\\[3pt]
    {\large Senior AI / LLM Engineer}\\[5pt]

    Mexico City, Mexico
    \enspace|\enspace
    \href{mailto:othon.gonzalezc@gmail.com}
    {othon.gonzalezc@gmail.com}
    \enspace|\enspace
    +52 55 3332 8699
    \\[2pt]
    \href{https://www.linkedin.com/in/thongc/}
    {linkedin.com/in/othongc}
    \enspace|\enspace
    \href{https://github.com/uthynauta}
    {github.com/uthynauta}
\end{center}

% ============================================================
% SUMMARY
% ============================================================

\section*{Professional Summary}

Senior AI / LLM Engineer with a PhD in Advanced Technology and experience
building production-oriented enterprise generative AI platforms, agentic
workflows, Retrieval-Augmented Generation systems, evaluation pipelines,
backend services, and LLM observability capabilities.

Strong background in Python, FastAPI, LangGraph, OpenTelemetry,
PostgreSQL/pgvector, Docker, Kubernetes, PyTorch, Transformers, and
Hugging Face. Experienced in designing systems that capture model usage,
cost, latency, tool execution, errors, execution paths, and agent outcomes
to improve debugging, evaluation, governance, reproducibility, and
model-routing decisions.

Previous experience in automotive perception, computer vision, NLP,
multimodal systems, and safety-critical engineering brings rigorous
experimentation, validation, and systems thinking to production AI
development.

% ============================================================
% CORE SKILLS
% ============================================================

\section*{Core Skills}

\textbf{Generative AI and Agents:}
Large Language Models, Agentic AI, LangGraph, tool calling, reusable agent
skills, Model Context Protocol (MCP), prompt orchestration, model routing,
Retrieval-Augmented Generation, retrievers, embeddings, vector search

\textbf{LLM Evaluation and Observability:}
OpenTelemetry, distributed tracing, token and cost tracking, latency analysis,
execution-path tracing, tool-call telemetry, error analysis, DeepEval,
evaluation datasets, agent outcome analysis, AI provenance

\textbf{Backend and Data:}
Python, FastAPI, PostgreSQL, pgvector, SQL, REST APIs, asynchronous workflows,
data pipelines, structured event models, JSON-based telemetry

\textbf{Infrastructure and Operations:}
Docker, Docker Compose, Kubernetes, Linux, Bash, Grafana, Prometheus, Loki,
containerized services, cloud-native development

\textbf{Machine Learning:}
PyTorch, Transformers, Hugging Face, TensorFlow, NLP, computer vision,
image captioning, multimodal learning, model evaluation

\textbf{Additional Engineering Experience:}
ADAS, radar perception, camera and LiDAR validation, multi-sensor systems,
safety-critical engineering, MATLAB, C++

\textbf{Languages:}
Spanish --- Native
\enspace|\enspace
English --- Professional working proficiency

% ============================================================
% EXPERIENCE
% ============================================================

\section*{Professional Experience}

\textbf{Teradata}
\hfill
\textit{Dic 2025 -- Jul 2026}

\textit{AI Engineer / Agentic AI Platform Engineer}

\begin{itemize}
    \item Developed core capabilities for an internal enterprise Agentic AI
    platform deployed to production for internal use, supporting agent
    development, execution, evaluation, observability, and reusable skill
    integration.

    \item Engineered LLM-based workflows using LangGraph, structured tool
    execution, reusable agent skills, MCP-based integrations, and
    Retrieval-Augmented Generation.

    \item Contributed to a scalable platform architecture initially supporting
    five production agents, designed to enable the onboarding, deployment,
    versioning, and operation of additional agents and reusable skills.

    \item Designed an OpenTelemetry-based observability architecture that
    standardized the capture of token usage, model cost, latency, errors,
    model metadata, tool calls, execution paths, and agent outcomes across
    multi-step workflows.

    \item Integrated telemetry and monitoring components using Grafana,
    Prometheus, Loki, and distributed tracing pipelines to support debugging,
    performance analysis, operational visibility, and AI FinOps.

    \item Built structured evaluation and analytics pipelines using
    PostgreSQL, pgvector, and DeepEval, enabling agent assessment, dataset
    generation, failure analysis, and performance-informed model-routing
    decisions.

    \item Defined telemetry schemas and event models for recording agent,
    graph, model, prompt, tool, and skill metadata while supporting
    redaction and controlled storage of sensitive information.

    \item Developed backend services and integration endpoints using Python
    and FastAPI for agent execution, platform services, and internal
    system interoperability.

    \item Developed containerized multi-service environments using Docker
    Compose and Kubernetes to support local development, integration testing,
    and production deployment.

    \item Collaborated on architecture decisions involving agent interfaces,
    skill discovery, model access, observability, evaluation, reproducibility,
    deployment patterns, and production readiness.
\end{itemize}

\textbf{Continental Autonomous Mobility}
\hfill
\textit{Jul 2022 -- Nov 2025}

\textit{Algorithms Test / Data Engineer --- ADAS}

\begin{itemize}
    \item Developed AI-assisted testing and data-processing tools for radar
    and multi-sensor perception systems, improving the efficiency and
    repeatability of model-validation workflows.

    \item Generated and curated simulated radar datasets in dSPACE for training and evaluating radar-virtualization models under controlled and reproducible scenarios.

    \item Built automated engineering workflows for test-data processing,
    validation activities, and repetitive analysis tasks using Python and
    workflow-automation tooling.

    \item Designed and executed validation strategies for radar, camera, and
    LiDAR perception pipelines, supporting system reliability in
    safety-critical automotive environments.

    \item Applied computer vision models and tooling including DETR, YOLO,
    Segment Anything, and CVAT in perception-data and validation workflows.

    \item Supported cross-functional collaboration among algorithm,
    validation, data, simulation, and systems-engineering teams.

    \item Contributed to a strategic cooperation agreement with the National
    Polytechnic Institute, enabling joint research and development
    opportunities.
\end{itemize}

\textbf{CentroGeo}
\hfill
\textit{Jul 2021 -- Jul 2023}

\textit{Postdoctoral Researcher}

\begin{itemize}
    \item Researched and developed NLP and multimodal deep-learning models
    for generative tasks, with emphasis on automatic image and video
    captioning.

    \item Architected and trained models based on recurrent neural networks,
    LSTMs, attention mechanisms, and Transformers for vision-language
    research problems.

    \item Designed experimental workflows for evaluating caption-generation
    systems and analyzing whether commonly used metrics reflected semantic
    and perceptual output quality.

    \item Used PyTorch, Transformers, and Hugging Face tools for model
    development, experimentation, training, and inference.

    \item Managed datasets and computational workflows for model training,
    inference, comparison, and reproducible experimentation.

    \item Evaluated generative-image technologies including DALL-E and
    Stable Diffusion for research-oriented applications.

    \item Co-authored peer-reviewed publications on video captioning,
    multimodal-model evaluation, and computer vision.
\end{itemize}

% ============================================================
% EARLIER LEADERSHIP EXPERIENCE
% ============================================================

\section*{Earlier Engineering and Leadership Experience}

\textbf{External Consultant}
\hfill
\textit{Feb 2016 -- Jun 2018}

\textit{Safety Manager}

\begin{itemize}
    \item Led operational-safety and compliance initiatives for a Mexican
    ground-handling services provider operating in a regulated aviation
    environment.

    \item Supported certification of the Toluca station under IS-BAH
    Safety Stages 1 and 2.

    \item Coordinated with Mexican authorities to resolve a critical permit
    renewal and prevent operational disruption.
\end{itemize}

\textbf{Aerom\'exico}
\hfill
\textit{Feb 2012 -- Oct 2013}

\textit{Regional Airports Chief and Processes Analyst}

\begin{itemize}
    \item Served as liaison between regional airport operations and corporate
    headquarters, coordinating process alignment and operational follow-up.

    \item Owned and improved a national KPI for baggage irregularities,
    applying data analysis to strengthen tracking and operational control.

    \item Supported ISAGO certification renewal and participated in
    regulatory and operational audits.
\end{itemize}

\textbf{ITESM, UVM and CentroGeo}
\hfill
\textit{2019 -- 2022}

\textit{Adjunct Professor}

\begin{itemize}
    \item Developed and taught undergraduate and graduate courses in
    statistics, physics, aerodynamics, engineering, and specialized
    aeronautical subjects.
\end{itemize}

\newpage

% ============================================================
% EDUCATION
% ============================================================

\section*{Education}

\textbf{PhD in Advanced Technology}
\hfill
\textit{CICATA Quer\'etaro, Instituto Polit\'ecnico Nacional, 2020}

Research focus: Computer Vision and Remote Sensing

Thesis: \textit{Precise Remote Measurement of Thermal Radiance Dynamics
Using Unmanned Aerial Vehicles}

Final grade: 9.0/10

\vspace{4pt}

\textbf{MSc in Advanced Technology}
\hfill
\textit{CICATA Quer\'etaro, Instituto Polit\'ecnico Nacional, 2016}

Research focus: Aerodynamic Optimization and Morphing Systems

Thesis: \textit{Data-Driven Morphable Wing for Tracking and Mapping
Using Small Unmanned Aerial Vehicles}

Final grade: 9.5/10
\enspace|\enspace
\textit{Cum Laude}

\vspace{4pt}

\textbf{BEng in Aeronautical Engineering}
\hfill
\textit{ESIME Ticom\'an, Instituto Polit\'ecnico Nacional, 2010}

Thesis: \textit{Software Tool to Determine the Capacity of Different
Runway Systems}

Final grade: 9.4/10
\enspace|\enspace
\textit{Cum Laude}

% ============================================================
% PUBLICATIONS
% ============================================================

\section*{Selected Publications}

\begin{itemize}
    \item Gonz\'alez-Ch\'avez, O., et al.
    ``Are Metrics Measuring What They Should? Evaluation of Image
    Captioning Task Metrics.''
    \textit{Signal Processing: Image Communication}, Vol. 120, 2024.

    \item Moctezuma, D., et al.
    ``Video Captioning: A Comparative Review.''
    \textit{Computer Vision and Image Understanding}, Vol. 231, 2023.

    \item Gonz\'alez-Ch\'avez, O., et al.
    ``Radiometric Calibration of Infrared Thermal Cameras.''
    \textit{IEEE Transactions on Instrumentation and Measurement}, 2019.

    \item Gonz\'alez-Ch\'avez, O., et al.
    ``Thermal Radiation Dynamics of Soil Surfaces with UAVs.''
    \textit{Mexican Conference on Pattern Recognition}, 2019.
\end{itemize}

\end{document}